\section{Limitations and Future Work}
\label{sec:limitations}

\textbf{Fixed gate set, no continuous rotation angles.} The six single-qubit presets
(Pauli~$X$/$Y$/$Z$, Hadamard, $S$, $T$) and the control/target multi-qubit mechanism are
enough to build every worked experiment in Section~\ref{sec:experiments}, but exclude an
arbitrary-angle rotation gate ($R_x(\theta)$, $R_y(\theta)$, $R_z(\theta)$), which a more
advanced course would eventually need. As with CircuitCraft's own preset-cycling
interaction~\cite{picos2026circuitcraft}, this design keeps the interaction model consistent
with the remainder of the base game (right-clicking through a fixed, small list rather than a
custom numeric-entry interface), at the cost of the continuous control a full gate set
provides; exposing a numeric-entry interface for an arbitrary rotation angle is a natural
extension.

\textbf{Hard cap of three qubits.} Section~\ref{sec:pedagogy} explains the pedagogical
rationale for this cap (keeping the full $2^n$-basis-state probability distribution
displayable at once), but it also means QuantumCraft cannot demonstrate any algorithm whose
interest specifically depends on a qubit count beyond three -- Grover's
algorithm~\cite{grover1996fast} and Shor's~\cite{shor1997polynomial} both require considerably
more qubits to exhibit their asymptotic advantage meaningfully. Raising this cap is
straightforward in the simulator itself (the statevector representation scales as $2^n$
regardless of the cap value) but would require redesigning the result heads-up display
(Figure~\ref{fig:hud}), which currently assumes every basis state's probability fits on
screen simultaneously.

\textbf{No mid-circuit measurement or classical feedback.} A circuit, once scanned, is always
run start to finish before any measurement occurs (Section~\ref{sec:quantum-simulator}); there
is no mechanism for a measurement partway through a circuit to classically control a later
gate, a pattern used by more advanced protocols such as quantum teleportation. Supporting this
would require the world-scan algorithm (Section~\ref{sec:architecture}) to interleave scanning
and execution rather than building the whole circuit before running any of it.

\textbf{No persisted circuit state, no saved presets.} As with CircuitCraft's
components~\cite{picos2026circuitcraft}, no block entity in this mod persists its
right-click-selected preset across a world save and reload, and a circuit itself is never
saved as an object -- only the placed blocks that describe it. This is a deliberate
simplification given that QuantumCraft's circuits are re-run from scratch on every probe use
regardless (Section~\ref{sec:quantum-simulator}), but it does mean a player must re-verify a
gate's cycled-to preset after any world reload, exactly as a CircuitCraft player must for a
resistor's resistance or a diode's forward-voltage family.

\textbf{Verification remains informal.} As noted in Section~\ref{sec:validation}, the nine
closed-form and statistical checks used during development have not yet been captured as a
persisted, automatically re-run regression suite, and the tool's pedagogical effectiveness
has to date been assessed only informally by its author, rather than through a controlled
classroom study. A natural next step is a pre/post evaluation using an established
quantum-computing concept inventory, comparing a cohort using QuantumCraft as a supplementary
laboratory activity against a cohort using a conventional notebook-based simulator, in order
to test whether the game-based, first-person construction experience described in
Section~\ref{sec:pedagogy} produces a measurable difference in conceptual understanding,
rather than in engagement alone.

\textbf{Platform dependency.} As with CircuitCraft, the mod is necessarily tied to a specific
version of Minecraft and of the Fabric mod loader, both of which are externally maintained and
continue to evolve, and ongoing maintenance to track future game updates remains an
unavoidable cost of building instructional software atop an actively developed commercial
platform.
